\SourceLeaf{206}%
% A544-P0206-D01
Beschränkung auf ganze Werthe der Coëfficienten weglassen will, da es dabei nicht auf absolute Werthe, sondern nur auf ihr Verhältniss unter einander ankommt, allezeit ganze Zahlen finden, die den Messungsresultaten so nahe kommen, wie man nur will.

% A544-P0206-D02
Schliesslich wollen wir noch dem oben angeführten \textsc{Seeber}'schen Lehrsatze seine geometrische Bedeutung unterlegen. Wenn ein Parallelepipedum so beschaffen ist, dass keine seiner zwölf Kanten (unter denen je vier einander gleich sind) grösser ist, weder als eine der zwölf Diagonalen von Seitenflächen (die paarweise gleich sind), noch als eine der vier Diagonalen des Parallelepipedum: so ist der mit $\surd2$ multiplicirte Rauminhalt desselben nicht kleiner, als der Rauminhalt eines aus denselben Kanten gebildeten rechtwinklichten Parallelepipedum.
\BookRule{.28\linewidth}

\clearpage
\SourceLeaf{207}%
% A544-P0207-D01
\renewcommand{\BookRunningTitle}{Handschriftlicher Nachlass}
\thispagestyle{empty}
\pdfbookmark[0]{Handschriftlicher Nachlass}{handschriftlicher-nachlass}
\vspace*{.25\textheight}
\begin{center}
{\Large HANDSCHRIFTLICHER}\par\vspace{2.8cm}
{\Huge NACHLASS.}
\end{center}
\vfill

\clearpage
\SourceLeaf{209}%
% A544-P0209-D01
\renewcommand{\BookRunningTitle}{Nachlass. Analysis residuorum}
\pdfbookmark[0]{Solutio congruentiae X^m-1≡0. Analysis residuorum}{analysis-residuorum-pars-prior}
\begin{center}
{\large SOLUTIO CONGRUENTIAE $X^m-1\equiv0$.}\par\medskip
ANALYSIS RESIDUORUM. CAPUT SEXTUM. PARS PRIOR.
\end{center}
\BookRule{.18\linewidth}
\renewcommand{\Article}[1]{\par\medskip\Needspace{4\baselineskip}\pdfbookmark[1]{Art. #1}{analysis-residuorum-art-#1}\centerline{\textbf{#1.}}\nobreak\smallskip}
\Article{237}
In Cap. \textsc{iii} docuimus, congruentiam $x^n\equiv1$, si pro modulo accipiatur numerus primus $p$, habere $\mu$ radices, quando $\mu$ est maxima communis mensura numerorum $n$ et $p-1$, hasque radices cum radicibus congr. $x^\mu\equiv1$ penitus convenire. Quamobrem eum casum considerare sufficit, ubi $n$ est pars aliquota numeri $p-1$. Quod autem non modo congruentiae $x^n-1\equiv0$ sed cuiusvis alius solutio pro modulis quibuscunque ex solutione pro modulis, qui sunt numeri primi, possit derivari, iam passim est ostensum infraque (Cap. \textsc{viii}) fusius docebitur.

% A544-P0209-D02
\Article{238}
Sed ne hic quidem subsistere opus est; namque eodem Capite \textsc{iii} exposuimus, congruentiae $x^n\equiv1$ solutionem a resolutione similium congruentiarum pendere $x^a\equiv1$, $x^b\equiv1$ etc., ubi $a,b$ etc. sunt numeri primi aut numerorum primorum potestates et $n$ productum ex his numeris. Si scilicet $A,B$ etc. sunt respective radices quaecunque congruentiarum $x^a\equiv1$, $x^b\equiv1$ etc., productum ex his $AB\ldots$ erit aliqua e radicibus congruentiae $x^n\equiv1$. Nostrae igitur investigationes ad solutionem congruentiae $x^n\equiv1\;(\mathrm{mod.}\,p)$ restringentur, quando $p$ est numerus primus, $n$ numerus primus aut numeri primi potestas, simulque pars aliquota numeri $p-1$.

\SourceLeaf{210}%
% A544-P0210-D01
\Article{239}
Porro ex Cap. \textsc{iii} constat, inter congruentiae $x^n\equiv1$ radices semper aliquas dari, per quarum potestates omnes ceterae exhiberi possunt. Ita si $r$ designet huiusmodi radicem (\emph{primitivam} supra diximus, quando $n=p-1$, hancque expressionem hic quamquam significatione latiori retinebimus) omnes congr. propos. radices erunt
\[
1,\ r,\ rr,\ r^3\quad\dotsfive\quad r^{n-1}
\]
Huiusmodi ergo radices omni studio sunt investigandae, quoniam his inventis ceterae sponte patebunt. Brevitatis gratia quamcunque ipsius $r$ potestatem per exponentem uncis inclusum designamus, ita ut $(0)$ denotet unitatem, $(1)$ radicem quamcunque primitivam congruentiae $x^n\equiv1$, $(2)$ ipsius $(1)$ quadratum etc.: ita ut haec series $(0),\ (1),\ (2),\ (3),\dotsfour(n-1)$ omnes radices amplectatur. Ceterum constat, $(k)$ semper fore talem radicem primitivam, quoties $k$ ad $n$ est primus; i.~e. nostro casu (ubi $n$ est numeri primi $t$ potestas $=t^\nu$), quoties $t$ ipsum $k$ non dividit. Manifesto vero signa $(1)$, $(2)$ etc. per se sunt indeterminata; sed simulac ipsi $(1)$ valor aliquis determinatus tribuitur, omnia cetera determinata fient.

% A544-P0210-D02
\Article{240}
Quoniam radices primitivas prae ceteris investigare propositum est, has a ceteris primum separare oportet. Quod fiet, si e serie $(0),(1),(2)\ldots(n-1)$ omnes terminos $(k)$ eiiciamus, ubi $k$ per $t$ dividitur; quodsi autem $n$ est numerus primus seu $\nu=1$, unicus $(0)$ erit abrogandus. Priusquam vero ad disquisitionem radicum superstitum progrediamur, lectorem sedulo admonemus exempla aliquot sibi conficere, ut omnia, quae sine his forsan generalius dicta viderentur, in concreto intueri possit. Nos aliquod apponimus; sed non ideo superfluum erit alia proprio Marte elaborare.

% A544-P0210-D03
Sit $p=29$; $n=7$ et septenae congruentiae $x^7\equiv1\;(\mathrm{mod.}\,29)$ radices erunt $1,7,16,20,23,24,25$. Quoniam $n$ est numerus primus, omnes hae radices praeter $1$ erunt primitivae; posito igitur $7=(1)$ signa haec significabunt:
\[
\begin{array}{ccccccc}
(0)&(1)&(2)&(3)&(4)&(5)&(6)\\
1&7&20&24&23&16&25
\end{array}
\]

